% from ega0-22-fr.tex
\oldpage[0\textsubscript{IV-1}]{208}

  \item[\rm(iii)] Soient $B=k[T_1,\ldots,T_r]$ une algèbre de polynômes,
  $\mathfrak J$ un idéal de $B$ tels que $C$ soit isomorphe à
  $B/\mathfrak J$. Soit $\mathfrak q$ un idéal premier de $B$ contenant
  $\mathfrak J$ et soit $\mathfrak p=\mathfrak q/\mathfrak J$. Les conditions
  de (i) sont alors encore équivalentes à la suivante («~critère jacobien de
  Zariski~»)~:
  \begin{enumerate}
    \item[e)] Il existe un système d'éléments
    $(f_i)_{1\leq i\leq m}$ de $\mathfrak J$, engendrant l'idéal
    $\mathfrak J_{\mathfrak q}$ dans $B_{\mathfrak q}$ et des
    $k_0$-dérivations $D_j$ de $B$ dans lui-même $(1\leq j\leq m)$ tels que
    $\det(D_j(f_i))\notin\mathfrak q$.
  \end{enumerate}
\end{enumerate}
\end{proposition}

\begin{proof}
On sait que $C$ est toujours isomorphe à une $k$-algèbre de la forme
$B/\mathfrak J$. Comme $B$ est une $k$-algèbre formellement lisse (19.3.3) et
que $k$, étant séparable sur $k_0$, est une $k_0$-algèbre formellement lisse
(19.6.1), $B$ est une $k_0$-algèbre formellement lisse (19.3.5, (ii)).
L'équivalence des conditions a), b), c) de (i) résulte donc de (22.6.4), et
celle de a), b), c) dans (ii) résulte de (22.6.6)~; l'équivalence de a) et d)
dans (i) résulte de (22.5.8)~; comme tout localisé de $C\otimes_{k_0}k'$ est
aussi un localisé de $C_{\mathfrak p}\otimes_{k_0}k'$ pour un idéal premier
$\mathfrak p\in\operatorname{Spec}(C)$ convenable, l'équivalence de d) et b) dans (ii)
résulte de l'équivalence de d) et a) dans (i). Enfin, comme
$\Omega^1_{B/k_0}$ est un $B$-module projectif, l'équivalence du critère de
Zariski e) et des autres conditions de (i) découle de (19.1.12) et de
(22.6.3), compte tenu de (20.4.8).
\end{proof}

Zariski s'intéressait en fait à un critère différentiel de régularité pour les
anneaux locaux $C_{\mathfrak p}$. Comme il revient au même de dire qu'un anneau
local noethérien contenant un corps est régulier ou est formellement lisse en
tant qu'algèbre sur son sous-corps premier (19.6.4), on obtient aussitôt un tel
critère en remplaçant $k_0$ par le sous-corps premier de $k$ dans (22.6.7)~; en
particulier, on obtient le résultat suivant, que nous retrouverons plus tard
(IV, 6.12.5) comme cas particulier de résultats plus généraux de Nagata~:

\begin{corollary}[22.6.8]
\label{0.22.6.8-fr}
\textnormal{(Zariski).}
Soit $C$ une algèbre de type fini sur un corps. Alors l'ensemble des
$\mathfrak p\in\operatorname{Spec}(C)$ tels que $C_{\mathfrak p}$ soit un anneau local
régulier est ouvert dans $\operatorname{Spec}(C)$.
\end{corollary}

\begin{remark}[22.6.9]
\label{0.22.6.9-fr}
\begin{enumerate}
  \item[\rm(i)] Les énoncés de (22.6.7) sont encore valables si au lieu de
  supposer $k$ séparable sur $k_0$, on suppose seulement qu'il est de
  \emph{multiplicité radicielle finie}, autrement dit (19.6.6) qu'il existe
  une extension radicielle finie $k'_0$ de $k_0$ telle que le corps résiduel de
  l'anneau artinien local $k\otimes_{k_0}k'_0$ soit une extension séparable
  $k'$ de $k'_0$. Il existe alors un $k'_0$-monomorphisme
  $k'\longrightarrow k\otimes_{k_0}k'_0$ qui, composé avec l'homomorphisme
  canonique $k\otimes_{k_0}k'_0\longrightarrow k'$, donne l'identité, puisque
  $k'$ est une $k'_0$-algèbre formellement lisse (19.6.1)~; on en conclut que
  $C\otimes_{k_0}k'_0$ est muni d'une structure de $k'$-algèbre de type fini,
  et comme $k'$ est séparable sur $k'_0$, on peut appliquer (22.6.7) à cette
  $k'$-algèbre~; on en conclut notre assertion en appliquant (19.4.7). Comme
  nous n'aurons pas à nous servir de cette généralisation, nous en laissons le
  détail au lecteur. Nous ignorons par contre si les résultats de (22.6.7) sont
  valables sans \emph{aucune} hypothèse sur l'extension $k$ de $k_0$.
  \item[\rm(ii)] Sous les hypothèses générales de (22.6.7), soit $(k_\alpha)$
  une famille filtrante décroissante de sous-corps de $k$ contenant $k_0$ et
  telle que
  $\bigcap_\alpha k_\alpha(k^p)=k_0(k^p)$ (où $p$ est l'exposant
  caractéristique de $k_0$). Utilisant une méthode de décompte des dimensions
  due à Nagata, on peut montrer que, dans le critère jacobien (22.6.7, (iii)),
  on peut se borner à des dérivations $D_j$ de $B$ qui soient des
  $k_\alpha$-dérivations pour un $\alpha$ convenable.

  